\documentclass[11pt]{article}
\usepackage[margin=1in]{geometry}
\usepackage{amsmath,amssymb,amsthm}
\usepackage{booktabs}
\usepackage[colorlinks=true,linkcolor=black,urlcolor=blue]{hyperref}

\newcommand{\RP}{\mathbb{R}\mathrm{P}}
\newcommand{\R}{\mathbb{R}}
\newcommand{\Cinf}{C_\infty}
\newcommand{\Hinf}{H_\infty}

\title{Counting the regions of the Clebsch cubic\\ minus its 27 lines}
\author{}
\date{}

\begin{document}
\maketitle
\vspace{-3em}

\section{The question}

Let $X=V(f_3)\subset\mathbb{P}^3$ be the Clebsch cubic surface and let $g$ be the
nonic with $V(g)\cap X$ equal to the $27$ lines. The routing function
$r=g/q^5$, with $q>0$ a shifted sum of squares, has
\[
  |r|\longrightarrow 0 \quad\text{on the $27$ lines and as } \|x\|\to\infty ,
\]
so \texttt{connected\_components} returns the connected components of
$X_{\mathrm{aff}}\smallsetminus V(g)$, where $X_{\mathrm{aff}}=X(\R)\cap\R^3$ is the
affine part of the surface in the chosen chart.

Writing $\Hinf$ for the plane at infinity and $\Cinf=X\cap\Hinf$, we have
$X_{\mathrm{aff}}=X(\R)\smallsetminus\Cinf$, so
\[
  X_{\mathrm{aff}}\smallsetminus V(g)
  \;=\; X(\R)\smallsetminus\bigl(\,\text{$27$ lines}\;\cup\;\Cinf\,\bigr).
\]
The count we want is therefore the number of \emph{faces} of the curve arrangement
$\mathcal{A}=\{\text{27 lines}\}\cup\Cinf$ on the closed surface $X(\R)$. The
arrangement is a CW decomposition, so Euler's formula gives
\begin{equation}\label{eq:euler}
  F \;=\; \chi\bigl(X(\R)\bigr) - V + E .
\end{equation}
Everything below is the computation of $V$ and $E$.

\section{The Euler characteristic of \texorpdfstring{$X(\R)$}{X(R)}}

All $27$ lines of the Clebsch surface are real. A smooth real cubic surface whose
$27$ lines are all real is $\RP^2$ blown up at six real points. A real blow-up
replaces a point by an $\RP^1$, and $\chi(\RP^1)=0$, so each blow-up changes the
Euler characteristic by $-1$. Hence
\begin{equation}\label{eq:chi}
  \chi\bigl(X(\R)\bigr) \;=\; \chi(\RP^2)-6 \;=\; 1-6 \;=\; -5 .
\end{equation}

\section{The arrangement of the 27 lines}

Every line on a smooth cubic surface meets exactly $10$ of the others, so the
number of intersecting pairs is
\[
  \tfrac{1}{2}\cdot 27\cdot 10 \;=\; 135 .
\]

\subsection*{A general cubic surface}

If no three lines are concurrent, those $135$ pairs give $135$ distinct points and
each line carries $10$ of them. A line of $X(\R)$ is a circle, and $k$ marked
points cut a circle into $k$ arcs, so
\[
  V=135,\qquad E=27\cdot 10=270,\qquad F=-5-135+270=\boxed{130},
\]
which is the number quoted for a general real cubic surface with all lines real.

\subsection*{The Clebsch surface}

The Clebsch is not general: it carries $10$ \emph{Eckardt points}, where three of
the $27$ lines are concurrent. (It is the unique cubic surface with ten of them.)
Each Eckardt point accounts for $\binom{3}{2}=3$ of the $135$ intersecting pairs
but contributes only one point, so
\begin{equation}\label{eq:V}
  V \;=\; 135-3\cdot 10+10 \;=\; 115 .
\end{equation}
Let $e(\ell)$ be the number of Eckardt points on a line $\ell$. Of the ten lines
meeting $\ell$, exactly $2e(\ell)$ of them meet it at Eckardt points, two at each
such point; the remaining $10-2e(\ell)$ meet it at distinct simple points. So
$\ell$ carries
\[
  e(\ell)+\bigl(10-2e(\ell)\bigr) \;=\; 10-e(\ell)
\]
marked points, hence that many arcs. Since each Eckardt point lies on three lines,
$\sum_\ell e(\ell)=3\cdot 10=30$, and
\begin{equation}\label{eq:E}
  E \;=\; \sum_{\ell}\bigl(10-e(\ell)\bigr) \;=\; 270-30 \;=\; 240 .
\end{equation}
Substituting \eqref{eq:chi}, \eqref{eq:V} and \eqref{eq:E} into \eqref{eq:euler},
\[
  F \;=\; -5-115+240 \;=\; \boxed{120}
\]
faces in $\RP^3$. The ten Eckardt points cost exactly ten regions relative to the
general surface.

\section{The curve at infinity}

The computation is affine, so $\Cinf$ must be added to the arrangement. It is a
genuine extra curve: restricting the degree-three part of the defining equation to
random real lines yields $1$ or $3$ real roots depending on the line, whereas three
real lines would give three every time. So $\Hinf$ is \emph{not} a tritangent plane
and $\Cinf$ is not a union of three of the $27$ lines.

Each line $\ell$ meets $\Hinf$ in a single point, which lies on $X$ and hence on
$\Cinf$: twenty-seven incidences, but numerically only
\[
  21 \text{ distinct points},\qquad\text{of which } 3 \text{ are already vertices.}
\]
Those three are Eckardt points lying at infinity; they carry $3\cdot 3=9$ of the
lines between them. Hence:

\paragraph{Vertices.} The new vertices are the points of $\ell\cap\Hinf$ not already
counted in \eqref{eq:V}:
\begin{equation}\label{eq:Vp}
  V' \;=\; 115 + (21-3) \;=\; 133 .
\end{equation}

\paragraph{Edges.} A line gains one further marked point unless its point at
infinity was already one of its marked points; the $9$ lines through the three
Eckardt points at infinity gain nothing, the other $18$ gain one. The curve
$\Cinf$ itself now carries the $21$ marked points and so contributes $21$ arcs.
\begin{equation}\label{eq:Ep}
  E' \;=\; \underbrace{240}_{\text{\eqref{eq:E}}} + \underbrace{18}_{\text{new marks}}
        + \underbrace{21}_{\Cinf} \;=\; 279 .
\end{equation}
(If some component of $\Cinf$ carried no marked point, it would need one dummy
vertex \emph{and} one dummy edge, leaving $V'-E'$ unchanged; the number of
components of $\Cinf$ therefore does not affect the answer.)

\paragraph{Faces.} By \eqref{eq:euler}, \eqref{eq:Vp} and \eqref{eq:Ep},
\[
  F' \;=\; -5 - 133 + 279 \;=\; \boxed{141}.
\]

\section{Why 141 is also the number of routing points}

Fix a face $C$ of the arrangement. On $C$ the function $|r|$ is positive, vanishes
on the whole of $\partial C$ --- the boundary consists of arcs of the $27$ lines,
where $g=0$, and of $\Cinf$, where $|r|\to 0$ because $\deg q^5=10>9=\deg g$ ---
and is proper. It therefore attains an interior maximum, which is a critical point
of $r|_X$, i.e.\ a routing point.

Morse theory gives $\chi(C)=\sum_p(-1)^{\operatorname{index}(p)}$, summed over the
critical points in $C$. Every face here is an open disc, so $\chi(C)=1$ and
generically $C$ contains exactly one critical point, of index $0$. Summing over
faces predicts
\[
  \textbf{141 routing points, every one of index } 0 .
\]
Because all indices vanish, the connectivity matrix of
\texttt{find\_connectivity\_matrix} is the identity: no gradient path is tracked,
and the number of components is simply the number of critical points found.

\section{Numerical confirmation}

The combinatorial input was computed from the surface itself, and the routing
points independently by the package.

\begin{center}
\begin{tabular}{lcc}
\toprule
Quantity & Predicted & Computed \\
\midrule
real lines                              & 27  & 27 \\
intersecting pairs                      & 135 & 135 \\
distinct intersection points $V$        & 115 & 115 \\
Eckardt points                          & 10  & 10 \\
lines with $8$ / $10$ marked points     & 15 / 12 & 15 / 12 \\
arcs $E$                                & 240 & 240 \\
\midrule
distinct points of $\ell\cap\Hinf$      & --- & 21 \\
of those, already vertices              & --- & 3 \\
lines gaining a marked point            & --- & 18 \\
\midrule
\textbf{faces} $F'$                     & \textbf{141} & --- \\
\textbf{routing points found}           & ---          & \textbf{141}, all of index $0$ \\
\bottomrule
\end{tabular}
\end{center}

The two counts are obtained independently, so their agreement is mutual
confirmation. Landing on exactly $141$ rather than more also confirms, after the
fact, the assumption that every face is a disc: an annular face would have
contributed an extra critical point of index $1$, and none was found.

\section{Caveats}

\begin{itemize}\itemsep2pt
\item The incidence counts ($135$ pairs, $10$ Eckardt points, and the $21/3/18$
  figures at infinity) are numerical, decided with a tolerance of $10^{-6}$ on
  normalised homogeneous coordinates.
\item $\Cinf$ is assumed smooth, i.e.\ $\Hinf$ is nowhere tangent to $X$. A
  tangency would change both $V'$ and $E'$.
\item The step from $141$ faces to $141$ critical points assumes each face is a
  disc carrying a single non-degenerate maximum of $|r|$; this is confirmed a
  posteriori by the computation.
\item $141$ counts faces in this particular affine chart. A different chart moves
  $\Cinf$ and changes the answer; the chart-independent number is the projective
  one, $120$.
\end{itemize}

\end{document}
